% !TeX spellcheck = en_us
\documentclass[a4paper,12pt,fleqn]{article}

\usepackage[left=1.50cm, right=1.50cm, top=3cm, bottom=2cm, 
headheight= 2.5cm, headsep= 4pt, footskip=25pt]{geometry}
\usepackage{fancyhdr}
\usepackage{lmodern}
\usepackage[table]{xcolor}
\usepackage{amssymb,amstext,amsmath}
\usepackage{eso-pic}
\usepackage{booktabs}
\usepackage{bm}
\usepackage{setspace}
\usepackage{graphicx}
\usepackage[colorlinks]{hyperref}
\setlength{\parindent}{0mm}
\setlength{\parskip}{2mm}
\usepackage{enumitem}
\usepackage{titlesec}
\usepackage{gensymb}
\usepackage{pgfgantt}
\usepackage{lscape}
\usepackage[absolute]{textpos}
\usepackage{comment}
\usepackage{cite}
\usepackage{tikz}
\usepackage{pgfplots}
\usetikzlibrary{shapes.geometric}
\usepackage{array}
\newcolumntype{P}[1]{>{\centering\arraybackslash}p{#1}}
\usepackage{longtable}
\usepackage[font={small}]{caption}
\usepackage[most]{tcolorbox}
\usepackage{tikz}
\usepackage{stmaryrd}
\usepackage{wrapfig}
%\usepackage[none]{hyphenat}
\usepackage{ctable}
\usepackage{ragged2e}
\usepackage[official]{eurosym}
\newcommand*\circled[1]{\tikz[baseline=(char.base)]{
		\node[shape=circle,draw,inner sep=1pt] (char) {#1};}}
\usepackage[symbol]{footmisc}
\renewcommand{\thefootnote}{\fnsymbol{footnote}}
\renewcommand{\thempfootnote}{\fnsymbol{mpfootnote}}
\renewcommand*\footnoterule{}

\usepackage{xcolor}
\definecolor{docColor}{cmyk}{1.00,1.00,0, 0.40}  % 0.90\% of black
\color{docColor}

\newtcolorbox{myhlbox}[2][]{colback=docColor!05, colframe=docColor!20, width=1.0\linewidth,
	boxrule=0pt, left=1em, right=1em, top=1em, bottom=1em,
	colbacktitle=docColor!20, coltitle=docColor,enhanced, attach boxed title to top left={xshift=2mm, yshift=-1mm},
	title=#2, #1}

\hypersetup{
	colorlinks=true,
	citecolor=blue,
	filecolor=magenta,
	linkcolor=blue,
	urlcolor=blue,
	pdftitle={SeminarAnnouncements},
	%bookmarks=true,,
}

\graphicspath{{./figs/}}

\definecolor{dgreen}{rgb}{0.0, 0.5, 0.0}

\allowdisplaybreaks

\newcommand*\mystrut[1]{\vrule width0pt height0pt depth#1\relax}
\renewcommand{\thesection}{\arabic{section}.}
\renewcommand{\thesubsection}{\arabic{section}.\arabic{subsection}}
\renewcommand{\thesubsubsection}{\arabic{section}.\arabic{subsection}.\arabic{subsubsection}}
\titleformat{\section}{\sffamily\large\bfseries}{\thesection}{1.0em}{}
\titleformat{\subsection}{\sffamily\large\bfseries}{\thesubsection}{1.0em}{}
\titleformat{\subsubsection}{\sffamily\normalsize\bfseries}{\thesubsubsection}{1.0em}{}
\titlespacing\section{0pt}{0pt plus 2pt minus 2pt}{0pt plus 2pt minus 2pt}
\renewcommand{\baselinestretch}{1.00}

\newcommand\CircledImage[1]{%
	\tikz\path[fill overzoom image={#1}, line width=0mm]circle[radius=0.5\linewidth];%
}

\input{format/defs}

\fancypagestyle{titlepage}{
	\fancyhf{}
	\insertLXheader
}

\fancypagestyle{plain}{%
	\fancyhf{}
	\renewcommand{\headrulewidth}{0pt}
	\renewcommand{\footrulewidth}{0pt}
	\insertnormalpage
}

\pagestyle{plain}

\newcommand\TitleBackground{%
	\AddToShipoutPictureBG*{
		\begin{tikzpicture}[remember picture,overlay]
			\node[opacity=0.20, anchor=center] 
			at (current page.center)
			{\includegraphics[height=\paperheight,keepaspectratio]{format/lms-bg.jpg}};
		\end{tikzpicture}
	}
}

\begin{document}
	\sffamily
	\thispagestyle{titlepage}
	\TitleBackground
	\vspace*{-3em}
	\begin{center}
		\huge \textbf{LMS Seminar}
	\end{center}
	\begin{center}
		\large 
		\textbf{A Mechanical description of Thermal conductivity and its Applications} 
	\end{center}
	\begin{center}
		\Large
		Anne Tanguy \\ 
		{\small  $^{\text{1}}$LaMCoS, CNRS, INSA Lyon~$^{\text{2}}$ONERA, University Paris-Saclay~$^{\text{3}}$LMS, Ecole Polytechnique, IP Paris} \\
		\vspace*{1em}
		{\large \textbf{Date, time, and venue} \\
		February 19, 2026 (2 -- 3 pm), Amphithéâtre Pole Mecanique}
	\end{center}

	\vspace*{-1em}
	
	\begin{myhlbox}{{Abstract}}
		The heat conductivity characterizes heat transfer which is a phenomenon of energy transport in matter, linking energy flux to temperature gradient. Two types of heat carriers have been already identified in solids: electrons and lattice vibrations, the latter also being described in terms of phonons in quantum mechanics. In this presentation, I will first remind the different descriptions of heat carriers and their connexion to thermal conductivity, within the kinetic theory and its link to quantum theory. We will show that eigenvibrations in solid disordered materials can take different forms, and that the phononic contribution to heat transfer can be simply described in terms of the propagation of mechanical wave packets. This mechanics-based approach not only allows us to obtain the non-trivial dependence of heat conductivity on temperature, also observed experimentally in careful experiments, but also to predict the role of nanostructural design in heat management in solids. Among the applications, we will discuss the possibility of manufacturing integrated circuits for heat management.
	\end{myhlbox}

\vspace*{1em}

	\begin{myhlbox}{{About the speaker}}
		\begin{wrapfigure}{r}{0.25\textwidth}
			\vspace*{-2em}
			\CircledImage{images/tanguy}
		\end{wrapfigure}

		Graduated from Ecole Polytechnique, Anne Tanguy began her professional carreer as an assistant professor of theoretical physics where her work focused on constitutive laws of glasses, then became a professor in mechanical engineering at INSA Lyon in 2015. In 2020, she has been appointed scientific director of the Materials and Structure division at ONERA. After spending a few months in the Ministry of Research as deputy director of the Research and Innovation Strategy department, she joined LMS and M4S as a professor at Ecole Polytechnique and dean of the bachelor.
		
	\end{myhlbox}
\end{document}